\expandafter\def\csname HistoricalNoteText491\endcsname{\emph{Ueber die Berechnung der Cometenbahnen.} \emph{Von} \textsc{Franz Güssmann}. \emph{Wien, bei Joh. Thom. Edlen von Trattnern}, 1803. 4to. 43 Seiten; mit einer Kupferplatte. Eigentlich hätte wohl eine so armselige mit Kuttenwitz ausgeschmückte Schmähschrift gar keine Erwähnung verdient, wenn nicht der Umstand, dass es gut sei, einen Menschen in seiner wahren Gestalt und Blösse darzustellen, welcher nur umherschleicht, um verdienten Männern bei nicht gehörig unterrichteten Personen zu schaden, Rücksicht verdiente. Von jeher hat dieser neidische und hämische Jesuit nichts als Schmähschriften, \emph{à la Valencia}, geschrieben; auf welche gewissenhafte und einsichtsvolle Art dieser heilige und gelehrte Mann dieses thut, davon hat man hier oben ein Pröbchen, dergleichen wir, wenn es der Mühe lohnte, noch mehrere aufstellen könnten.}
\clearpage
\fancyhead[L]{\small GÜSSMANN. COMETENBAHNEN.}
\hypertarget{native-491-lower}{}\label{native-491-lower}
\begin{center}
\small Monatliche Correspondenz. Band XV. Seite 452..460. 1807 Mai.\\[3pt]
\rule{.72\linewidth}{.3pt}
\end{center}
\begin{center}Auszug aus einem Schreiben.\end{center}
\begin{flushright}--- --- --- den 10. März 1807.\end{flushright}

--- --- --- Hierbei habe ich die Ehre, Ihnen die verlangte Brochüre\HistoricalNote{491} von \textsc{Güssmann} zu übersen%
\SourceLeaf{492}%
den. Einer ausführlichen und detaillirten Beantwortung ist sie wohl nicht werth; diese könnte nur dadurch interessant werden, dass man zwischen den Methoden von \textsc{Olbers} und \textsc{Boscovich} eine Parallele zöge und die Vorzüge, welche die erstere vor der letztern in Ansehung der grössern Bequemlichkeit hat, in's Licht stellte. Man kann von diesen Vorzügen, \textsc{Boscovich}'s Ruhm unbeschadet, sehr überzeugt sein, allein von diesem Ruhme kommt natürlich seinem Ordensbruder nichts zu gut. Die Würdigung von des letztern Brochüre, die keineswegs dazu geeignet ist, bei dieser Untersuchung mit in Betrachtung zu kommen, wird also wohl schicklicher davon getrennt bleiben.

Auf die ungezogenen, armseligen, zur Sache gar nicht gehörenden Einfälle sich nur mit einem Worte einzulassen, würde Ihrer und des Dr. \textsc{Olbers} sehr unwürdig sein. Ich begnüge mich, nur die Hauptpunkte heraus zu heben, womit er Unkundigen etwas gegen Dr. \textsc{Olbers} Methode gesagt zu haben scheinen könnte.

\textsc{Güssmann}'s vornehmster Angriff ist gegen Dr. \textsc{Olbers} Formel gerichtet, womit das Verhältniss der curtirten Abstände in der ersten und dritten Beobachtung bestimmt wird. \textsc{Güssmann} gibt dann in \emph{seiner Manier} eine Ableitung, worin er denn eine Menge Ungereimtheiten und Fehler findet; allein eine solche Ableitung wird Dr. \textsc{Olbers} nicht für die seinige erkennen, und jeder, der des letztern klaren Vortrag mit \textsc{Güssmann}'s wunderlicher Confusion der Begriffe vergleicht, wird dazwischen gar keine Uebereinstimmung finden. Ich weiss nicht, ob ich diese durchgängige Entstellung des \textsc{Olbers}'schen Vortrags für \emph{absichtlich} halten, oder glauben soll, dass \textsc{Güssmann} ihn selbst nicht verstanden habe. Gleich im 1., 2. und nachher im 21. §. wird \textsc{Olbers} die Voraussetzung \emph{angedichtet}, dass die Bewegung des Cometen als geradlinig angesehen werde. Zwischen Dr. \textsc{Olbers} einziger Voraussetzung, dass die Chorden der zwischen der ersten und letzten Beobachtung durchlaufenen Bogen der Cometen- und Erdbahn von den mittlern \emph{Radiis vectoribus} im Verhältnisse der Zeiten geschnitten werden, und der, welche \textsc{Güssmann} ihm \emph{aufdringen} will, dass der Comet und die Erde sich in der Chorde selbst bewegen, ist ein so grosser Unterschied, dass eine Gedankenlosigkeit oder — eine Unverschämtheit ohne Gleichen dazu gehört, ihn unwissend oder wissend zu übersehen. \textsc{Olbers} Ableitung des oben gedachten Verhältnisses ($M$) ist so lichtvoll, als irgend ein Leser, dem seine Schrift bestimmt ist, erwarten kann. Es seien $k$ und $i$ die äussern geocentrischen Oerter des Cometen auf der Himmelskugel; $h$ der Ort, wo der Comet der Erde erscheinen würde, wenn jener und diese in denjenigen Punkten der zwischen den äussern Oertern gezogenen Chorden stünden, wo diese von den mittlern \emph{Radiis vectoribus} geschnitten werden; $h^*$ der wirkliche geocentrische Ort in der mittlern Beobachtung auf der Himmelskugel; $g$ der mittlere Ort der Sonne: Dann liegen einerseits $g$, $h$, $h^*$ in einem grössten Kreise, andererseits auch $k$, $h$, $i$. Das erstere, weil offenbar jene fingirte Gesichtslinie in einer Ebene mit der Sonne und den mittlern Oertern der Erde und des Cometen liegt. Das andere ist eine leichte — aber von \textsc{Olbers} mit Recht seinen Lesern zur Entwickelung überlassene — Folge der Voraussetzung, dass die Stücke, worin die Chorden der Erdenbahn und Planetenbahn von den mittlern \emph{Radiis vectoribus} geschnitten werden, einander proportional sind. Hierauf gründet nun \textsc{Olbers} eine leichte Methode, die Lage von $h$ auf der Himmelskugel zu bestimmen; dann wird hiermit das Verhältniss der curtirten Abstände in den äussern Beobachtungen $M$ bestimmt; und endlich zeigt \textsc{Olbers}, wie man auch der vorläufigen Bestimmung von
\SourceLeaf{493}
$h$ ganz überhoben sein und unmittelbar $M$ aus den Datis des Problems berechnen könne. Dieses letztere hätte analytisch durch die schon vorher gefundenen Formeln geschehen können; \textsc{Olbers} hat aber für anschaulicher gehalten, es mit Hülfe einer Projection zu entwickeln. Man vergleiche nun diesen lichtvollen Vortrag mit \textsc{Güssmann}'s abenteuerlicher Confusion. Im 6. §. und in der Fig. II. bezeichnet er, wie ich, die äussern Oerter durch $k$, $i$, aber durch $h$ bezeichnet er erst den \emph{wahren Ort}: und behauptet, \textsc{Olbers} setze voraus, $h$ läge mit $k$ und $i$ in einerlei grösstem Kreise; dies ist abermals eine unverschämte Unwahrheit, denn das hat \textsc{Olbers} nirgends gesagt; aber gleich nachher nimmt er $h$ für den oben so bezeichneten fingirten Ort. Eine solche Verwechslung bringt ihn nun ganz in Verwirrung. Er meint, durch \textsc{Olbers} Verfahren müsse man für $c''$, d.\,i. die Länge des Punktes $h$, immer die Länge des Punktes $h^*$, die durch $a''$ bezeichnet ist, wieder erhalten; zu welchem sonderbaren Urtheile ihn nur die \emph{willkührliche} Verwechslung der Punkte $h$ und $h^*$, und die ganz ungegründete Voraussetzung, dass die drei geocentrischen Oerter in einem grössten Kreise liegen, gebracht hat. Er kann überhaupt die Möglichkeit gar nicht begreifen, die Reduction $c''-a''$ aus den Beobachtungen unmittelbar zu berechnen, weil dieselbe von zwei noch unbekannten Elementen abhänge, nemlich von den Abständen des Cometen und der Erde von der Sonne. Dies ist allerdings richtig; aber die Sache erklärt sich leicht daraus, dass von \emph{eben} diesen Elementen auch die Ausweichung des mittlern Ortes von dem durch die äussern gelegten grössten Kreise abhängig ist, und dass gerade diese Abweichung eigentlich dient, um $c''-a''$ zu bestimmen; dass jene Ausweichung gewöhnlich sehr klein ist, thut hier nichts zur Sache, denn $c''-a''$ ist damit von einerlei Ordnung. Die Vorstellung, die \textsc{Güssmann} im 7. §. von \textsc{Olbers} Gebrauch der orthographischen Projection gibt, ist wieder ganz unrichtig. In seiner Fig. III haben die Punkte $b$, $i$, $r$ gar nichts mit unserer Sache zu thun, sondern blos $z$ und $n$. Die von $z$ nach $n$ gezogene Linie bezieht sich auf den erwähnten fingirten Ort, aber nicht die von $b$ nach $r$ gezogene. Die Neigung der Projection der Gesichtslinie $z\,n$ ist es, die \textsc{Olbers} durch $b''$ bezeichnet und für die also nothwendig
$\operatorname{tang} b'' = \dfrac{\operatorname{tang} \gamma''}{\sin(A''-c'')}$
wird. Dass nun aber $\dfrac{\operatorname{tang} \gamma''}{\sin(A''-c'')} = \dfrac{\operatorname{tang} \ObservedBeta''}{\sin(A''-a'')}$, folgt unmittelbar daraus, dass die Projectionen von $z\,n$ und von $\mathfrak{S}\,m$ zusammenfallen. Wenn \textsc{Güssmann} nachher im 9. §. meint, \textsc{Olbers} Formel müsse eigentlich $\dfrac{t''\sin(a''-a')}{t'\sin(a'''-a'')}$ sein, so ist dies immer die Folge seiner unbegreiflichen Vermengung der Punkte $h$ und $h^*$; wenn er aber die Verschiedenheit der \textsc{Olbers}'schen Formel von diesem Ausdrucke daraus erklärt, dass \textsc{Olbers} alle drei Breiten $\ObservedBeta'$, $\ObservedBeta''$, $\ObservedBeta'''$ aufnimmt, die alle drei zusammen für eine gerade Linie und zwar auf bestimmte Punkte derselben aus bestimmten Punkten der Erdchorde nicht passen können, so urtheilt \textsc{Güssmann} hier völlig richtig und gibt dadurch, ohne es zu wissen, den Grund an, warum \textsc{Olbers} Formel richtig, hingegen der Ausdruck $\dfrac{t''\sin(a''-a')}{t'\sin(a'''-a'')}$, welcher sich auf die Voraussetzung der geradlinigen Bewegung gründet, unrichtig ist. Denn in der That lässt sich zeigen — wozu aber nicht hier der Ort ist, dass diese letzte Formel nur dann eine brauchbare Annäherung geben kann, wenn die Distanzen der Erde und des Cometen von der Sonne nicht sehr ungleich sind. Dieser Umstand findet zufälligerweise bei \textsc{Olbers} beiden Beispielen Statt und blos daher kommt es, dass die Werthe von $M$ nach dem erwähnten Ausdrucke berechnet, von den nach der \textsc{Olbers}'schen Formel berechneten nicht mehr verschieden sind, als \textsc{Güssmann} im 9. §. findet. Eben diesem Umstand ist es zuzuschreiben, dass \textsc{Güssmann} in seinem Exempel §. 24. nach jenem Ausdrucke noch erträglich richtige Resultate findet. Die trivialen Bemerkungen, die \textsc{Güssmann} im 11. §. macht, sind gewiss auch dem Dr. \textsc{Olbers} nicht entgangen. Dass $M$ eben so sicher bei kleinen als bei grossen Werthen von $a'''-c''$ und $c''-a'$ durch $\dfrac{t''\sin(c''-a')}{t'\sin(a'''-c'')}$ bestimmt werde, wenn man $c''$ in vollkommener Schärfe hat, ist allerdings wahr; aber bei der begränzten Genauigkeit aller auf die Sinustafeln gegründeten Rechnungen bleibt bei dem berechneten Werthe von $c''$ immer eine Ungewissheit übrig; \textsc{Olbers} wollte
\SourceLeaf{494}
blos sagen, dass diese Ungewissheit, möchte sie auch unter $0''1$ sein, in Fällen, wo $a'''\!-c''$, $c''\!-a'$ sehr klein sind, den Werth von $M'$ entstellen könnte. Eben so ist \textsc{Güssmann}'s zweite Bemerkung, dass die Vernachlässigung des Unterschiedes zwischen $c''$ und $a''$, den \textsc{Güssmann} mit $\delta$ bezeichnet, auch da, wo die Zwischenzeiten noch so klein sind, theoretisch genommen, nicht verstattet werden dürfe, ganz recht; aber unrecht hat er, wenn er glaubt, dass $\delta$ ein um so beträchtlicherer Theil von $a''\!-a'$ sei, je kleiner die Zwischenzeit ist; dies verhält sich gerade umgekehrt, denn $a''\!-a'$ ist eine Grösse der ersten, $\delta$ aber eine der zweiten Ordnung. Dass aber dessenungeachtet auch für sehr kleine Zwischenzeiten die Vernachlässigung von $\delta$ einen bedeutenden Einfluss auf die Resultate haben könne, wusste \textsc{Olbers} sehr wohl, und wenn er diese Vernachlässigung dennoch in diesem Falle erlaubte, so geschah es, weil dieser Einfluss abstrahirt von seiner absoluten Grösse in Vergleichung gegen die aus den unvermeidlichen Beobachtungsfehlern entspringenden Unrichtigkeiten um so kleiner ist, je kleiner die Zwischenzeiten sind. In solchen Fällen, wo die Unsicherheit der Beobachtungen doch einmal nur eine rohe Bestimmung der Bahn erlaubt, würde man sehr unrecht thun, durch ängstliche Rücksicht auf solche Umstände, deren Einfluss viel geringer ist, als jene Unsicherheit, sich die Arbeit zu erschweren und dadurch eine Genauigkeit affectiren, die die Umstände nicht erlauben. In wie fern es übrigens in solchen Fällen zweckmässiger sein kann, die Rechnung lieber ganz zu unterlassen, als unsichere Resultate zu entwickeln, bleibt immer den verständigen Astronomen aus den speciellen Umständen selbst zu beurtheilen überlassen.

Hiermit schliesse ich meine Bemerkungen; es genügt mir, über die von \textsc{Güssmann} in theoretischen Rücksichten gemachten Angriffe gesprochen zu haben. Wer die Kürze und Bequemlichkeit der \textsc{Olbers}'schen Methode mit den von \textsc{Güssmann} nachher angestellten weitschweifigen und doch nur zu grob genäherten Resultaten führenden Rechnungen, zu vergleichen Lust hat, wird über diesen Punkt auch leicht ein Urtheil fällen. Die Sache spricht für sich; dass \textsc{Güssmann} \S.\,16. Herr \textsc{von Pacassi} wegen der nach \textsc{Boscovich}'s Methode berechneten aber ganz fehlerhaften Bahn des Cometen von 1779 damit rechtfertigt, dass es sehr leicht sei, in einer sehr \emph{weitläuftigen} Rechnung einen Fehler zu begehen, der hernach ein unrichtiges Resultat verursache --- dies ist doch ein sprechender Beweis für den Vorzug, den Dr. \textsc{Olbers} Methode in dieser Rücksicht unstreitig hat. Es sind bekannte Thatsachen, dass nicht einer, sondern viele, auch nicht Mathematiker von Profession, nicht solche, die auf den Namen von vorzüglich geübten Rechnern darum Anspruch machten, sich mit Leichtigkeit die \textsc{Olbers}'sche Methode zu eigen gemacht und Cometenbahnen sicher und richtig darnach berechnet haben.
